\section{Game Logic and Server Authority}
\label{sec:authority}

\begin{guidance}[Game logic and server authority (15 marks)]
Show: correct \textbf{move validation}; rejection of \textbf{illegal,
out-of-turn and out-of-state} moves with specific errors; reliable
\textbf{win and draw detection}; and a \textbf{client-untrusted} design. Most
importantly, give an argument that a \textbf{rejected request cannot change the
server state}, based on how your code is structured (e.g. every check runs
before the first write), and point to the tests that check it. The brief calls
this a ``proof''; make it a step-by-step argument a reader can verify against
your code.
\end{guidance}

\subsection{Every check returns before the first write}
Listing~\ref{lst:move} is \texttt{handle\_move}. It runs the checks in the
order of Section~\ref{sec:protocol}: game state, syntax, range, turn, cell. Each
failing check sends one NACK and returns. The parser uses only local variables,
and saturates the number at 1000 so that a long digit string cannot overflow;
an earlier version used \texttt{atoi}, which turned \texttt{MOVE:abc} into a
move on cell 0.

\begin{lstlisting}[float=tp, language=C, caption={Move validation (\texttt{main.c}, trimmed).}, label=lst:move]
void handle_move(const char *p) {
    if (game_state != STATE_IN_PROGRESS)
        { send_str("NACK:NOGAME\n"); return; }
    int neg = 0, nd = 0, idx = 0; char mark = 0;
    if (*p == '-') { neg = 1; p++; }
    while (*p >= '0' && *p <= '9') {
        if (idx < 1000) idx = idx*10 + (*p - '0');
        p++; nd++;
    }
    if (nd == 0) { send_str("NACK:BADCMD\n"); return; }
    if (*p == ':') {
        mark = p[1];
        if ((mark != 'X' && mark != 'O') || p[2] != '\0')
            { send_str("NACK:BADCMD\n"); return; }
    } else if (*p != '\0')
        { send_str("NACK:BADCMD\n"); return; }
    if (neg || idx > 8)
        { send_str("NACK:RANGE\n"); return; }
    if (mark && mark != turn)
        { send_str("NACK:NOTURN\n"); return; }
    if (board[idx] != '.')
        { send_str("NACK:OCCUPIED\n"); return; }
    int r = apply_move(idx);  /* first write to the board */
    if (r == MV_CONTINUE && ai_mode) r = ai_reply();
    send_move_result(r);      /* WIN, DRAW or OK:BOARD */
}
\end{lstlisting}

\subsection{A rejected request cannot change the game state}
The game state is six variables: \texttt{board}, \texttt{turn},
\texttt{game\_state}, \texttt{connected}, \texttt{ai\_mode} and the move history.
The argument has five steps, and each step can be checked against \texttt{main.c}.
\begin{enumerate}\setlength\itemsep{0pt}
\item Only seven pieces of code write those variables: \texttt{game\_new()},
  \texttt{apply\_move()}, \texttt{handle\_undo()}, and the \texttt{CONNECT},
  \texttt{DISCONNECT}, \texttt{NEW} and \texttt{RESIGN} branches of
  \texttt{process\_command}. Every other path calls only \texttt{send\_str}.
\item \emph{Framing and checksum errors.} \texttt{process\_rx} calls
  \texttt{process\_command} only for a complete frame with \texttt{frame\_err}
  equal to zero and a checksum that is absent or correct. The oversize, garbage,
  fault, timeout and \texttt{CHECKSUM} branches send one NACK and reset
  \texttt{line\_len} and \texttt{frame\_err}, which belong to the framer, not to
  the game.
\item \emph{Unknown command or no session.} \texttt{process\_command} returns
  with \texttt{NACK:BADCMD} or \texttt{NACK:NOCONN} before it reaches any
  handler.
\item \emph{Handlers.} \texttt{NEW} returns with \texttt{NACK:BADCMD} before it
  sets \texttt{ai\_mode} or calls \texttt{game\_new()} when a game is running.
  \texttt{STATE} and \texttt{HISTORY} only read. \texttt{RESIGN} returns with
  \texttt{NACK:NOGAME} before it writes \texttt{game\_state}. \texttt{handle\_undo}
  checks the game state and the history length before it removes anything.
\item \emph{\texttt{MOVE}.} In Listing~\ref{lst:move} the call
  \texttt{apply\_move(idx)} is the first write, and every NACK path returns
  before it.
\end{enumerate}
Hence every path that ends in a NACK has written no game variable. This is an
argument from reading the code, not a formal proof, and the tests below sample
it. The button is not a request over the link: it is a local override that
resets the game when a session is open, and is documented as such.
\texttt{test\_occupied\_move\_rejected\_and\_board\_unchanged},
\texttt{test\_malformed\_move\_does\_not\_touch\_state},
\texttt{test\_new\_during\_game\_rejected\_state\_unchanged},
\texttt{test\_out\_of\_turn\_rejected},
\texttt{test\_wrong\_checksum\_rejected\_and\_nothing\_executed} and
\texttt{test\_undo\_all\_the\_way\_back\_then\_nothing\_to\_undo} each send a bad
request and then read \texttt{STATE} to compare the board. The oversize and
garbage tests check that a bad frame does not run the command at its end.

\subsection{The server alternates the turn and can check a named player}
The variable \texttt{turn} starts as X and flips only after a legal move that
neither wins nor draws. A plain \texttt{MOVE:<i>} is played by the side to move.
A \texttt{MOVE:<i>:<m>} with the wrong mark gets \texttt{NACK:NOTURN} before the
cell is examined, and \texttt{test\_out\_of\_turn\_rejected} checks that the board
is unchanged. In AI mode the human is always X, so a mark \texttt{O} on a
\texttt{MOVE} is also refused with \texttt{NACK:NOTURN}.

\subsection{The AI cannot make an illegal move}
\texttt{ai\_choose} considers only empty cells on a copy of the board, and
\texttt{ai\_reply} checks \texttt{board[idx] == '.'} again before it calls
\texttt{apply\_move}, the same function a human \texttt{MOVE} ends in. A legal AI
move is therefore true by construction, not by testing alone. Two checks back it
up. On a PC build of the same \texttt{main.c} I enumerated every possible human
strategy against the AI: 521 games, the human won none, the AI won 368 and 153
were draws.
\texttt{test\_ai\_never\_makes\_an\_illegal\_move\_and\_never\_loses} plays 20
seeded random games on the board and checks that each reply differs from the
previous board by exactly the human's X and one new O.

\subsection{Undo restores the board and the turn from the history}
The server stores the cell of every move. \texttt{UNDO} clears the last cell (the
last two in AI mode, so it is the human's turn again) and sets the turn from the
parity of the history length, so the turn can never disagree with the board. It
is refused after game over, because the result has already been sent, and with an
empty history it returns \texttt{NACK:NOUNDO}. \texttt{test\_undo\_restores\_board\_and\_turn}
and \texttt{test\_undo\_in\_ai\_mode\_removes\_both\_moves} check both modes.

\subsection{Win is checked before draw, for the mover only}
After placing a mark the server tests the eight lines (three rows, three columns,
two diagonals) for the mover's mark only, because only that mark can have just
completed a line. The win test runs before the draw test, so a ninth move that
completes a line is a win and not a draw. Either result sets
\texttt{GAME\_OVER}, after which \texttt{MOVE} returns \texttt{NACK:NOGAME}.
\texttt{test\_scripted\_win\_x}, \texttt{test\_scripted\_win\_o} and
\texttt{test\_scripted\_draw} cover these results. No test yet covers a win on the
ninth move, and Section~\ref{sec:conclusion} lists this as a limitation.

\subsection{The client holds no rules}
The GUI sends \texttt{MOVE:<i>} for any click, including a click on an occupied
cell, and shows the server's reason. If the client also checked legality, two
rule sets would have to stay identical, and they would disagree whenever the
button started a game that the client did not see. My own first GUI inferred
the final mark after a \texttt{WIN} from the move count. I removed that guess,
and the GUI now asks for \texttt{STATE}.
